\section{Abstract}
  Diese Masterarbeit beschäftigt sich mit der Entwicklung einer kombinierten Regelungsstrategie für ein autonom fahrendes,
  selbstbalancierendes Fahrrad, basierend auf den Ergebnissen dreier vorangegangener Forschungsarbeiten. Ziel dieser Arbeit 
  ist die erfolgreiche Integration einer PID-Balanceregelung mit einer MPC-basierten Fahrtrichtungsregelung, die mithilfe 
  eines YOLO-Modells zur visuellen Spurdetektion ergänzt wird. Zur Validierung und Optimierung der kombinierten Regelung 
  wurde die Simulationsumgebung Webots verwendet, wodurch die Robustheit und Funktionalität der entwickelten Regelungsalgorithmen
  realitätsnah überprüft werden konnten. Im Rahmen der Arbeit wurde eine zweistufige Regelungsarchitektur umgesetzt, welche die präzise 
  Balancierung des Fahrrads mittels eines PID-Reglers und die zuverlässige Spurverfolgung sowie Navigation mithilfe eines MPC-Reglers sicherstellt. 
  Die gewonnenen Simulationsergebnisse bestätigen die Machbarkeit und Effizienz des entwickelten Ansatzes und legen damit eine solide Grundlage für die zukünftige
  Implementierung der Regelung in einem realen autonomen Fahrrad.
\newpage